\chapter{Исследовательская часть}

\section{Замеры реального времени}

% Сравнить реальное время, затрачиваемое на обработку N \in {25, 50, 75, 100, 125} заявок последовательно и конвейером, основанным на k нативных потоках, не считая диспетчера.

Замеры проводились на ноутбуке Redmibook Pro 14 2022:
\begin{itemize}
	\item процессор Intel i7-12650H (архитектура x86-64, 10 ядер по 2.3 Ггц);
	\item оперативная память 16 Гб (DDR4, 5200 МГц);
	\item операционная система Windows 11 Home 25H2.
\end{itemize}
Для замеров использовалась функция now() из библиотеки chrono~\cite[с.~705]{lit1}.

Значение k=16 (получено в лабораторной работе №4). 

\begin{figure}[!htbp]
	\centering
	\begin{tikzpicture}
		\selectcolormodel{gray}
		\begin{axis}[
				axis lines=left, grid=both, minor grid style={gray!25}, major grid style={gray!50},
				legend style={font=\scriptsize, legend pos=north west}, ymin=0.002,
				xlabel=N, ylabel={Время, с},
			]
			\addplot+[smooth, mark=none] table [x index=0, y index=1] {../code/restm.txt};
			\addlegendentry{при k потоках}
			\addplot+[smooth, mark=none, dashed] table [x index=0, y index=2] {../code/restm.txt};
			\addlegendentry{ при 1 потоке}
		\end{axis}
	\end{tikzpicture}
	\caption{График зависимости времени от N заявок.}
\end{figure}

\section*{Вывод}

В данном разделе мной был проведен сравнительный анализ времени выполенения при различном количестве нативных потоков. 

\clearpage
